\documentclass[11pt,a4paper]{article}
\usepackage[utf8]{inputenc}
\usepackage[T1]{fontenc}
\usepackage[french]{babel}
\usepackage[top=2cm,bottom=2.5cm,left=2.5cm,right=2cm]{geometry}
\usepackage{enumitem}

\begin{document}

\begin{center}
{\large\bfseries Préparation au concours d'entrée à l'IFORD -- Yaoundé}\\[2mm]
{\Large\bfseries Épreuve de CULTURE GÉNÉRALE (commune Types A et B)}\\[2mm]
Banque de sujets d'entraînement \quad --\quad Durée : 4 heures\\[1mm]
\emph{Sujets rédigés pour l'entraînement, dans l'esprit de l'épreuve. Ce ne sont pas des sujets officiels.}
\end{center}
\hrule\medskip

\section*{Méthode}
L'épreuve consiste à traiter \textbf{un sujet au choix} (dissertation ou commentaire).
Le jury attend :
\begin{itemize}
  \item une \textbf{introduction} : amorce, définition des termes, problématique, annonce du plan ;
  \item un \textbf{développement} en deux ou trois parties équilibrées, argumentées et illustrées
  (chiffres de population, exemples africains, politiques publiques, conférences internationales) ;
  \item une \textbf{conclusion} : bilan et ouverture.
\end{itemize}
Repères utiles : Conférence internationale sur la population et le développement (Le Caire, 1994),
Objectifs de développement durable (2015--2030), dividende démographique, Agenda 2063 de l'Union africaine,
transition démographique, recensements et enquêtes démographiques et de santé (EDS).

\section*{Sujets}
\begin{enumerate}[label=\textbf{Sujet \arabic*.}, leftmargin=*]
  \item « La croissance démographique est-elle un frein ou un moteur du développement en Afrique subsaharienne ? »
  \item « Capturer le dividende démographique : illusion ou opportunité pour les pays africains ? »
  \item L'urbanisation rapide des villes africaines : défis et perspectives.
  \item « Investir dans la jeunesse, c'est investir dans l'avenir. » Commentez.
  \item L'autonomisation des femmes est-elle la clé de la maîtrise de la fécondité ?
  \item Les migrations internationales : une chance ou une menace pour les pays de départ ?
  \item « Sans données fiables, pas de politiques publiques efficaces. » Discutez cette affirmation à partir de l'exemple des recensements et de l'état civil.
  \item Changement climatique et mouvements de population en Afrique.
  \item Le numérique peut-il réduire les inégalités d'accès à l'éducation et à la santé ?
  \item Le vieillissement de la population concerne-t-il déjà l'Afrique ?
\end{enumerate}

\newpage
\section*{Corrigé indicatif -- Sujet 1 (plan détaillé)}

\paragraph{Introduction.}
L'Afrique subsaharienne connaît la croissance démographique la plus rapide au monde
(fécondité encore élevée, mortalité en baisse : transition démographique inachevée).
\emph{Définitions} : croissance démographique (accroissement naturel + solde migratoire) ;
développement (amélioration durable du bien-être : revenu, santé, éducation).
\emph{Problématique} : la dynamique démographique entrave-t-elle le développement, ou peut-elle,
sous conditions, en devenir un levier ?

\paragraph{I. Une croissance qui pèse sur le développement (thèse « néo-malthusienne »).}
\begin{enumerate}[label=\Alph*.]
  \item Pression sur les services sociaux : écoles, centres de santé, logements insuffisants face à la demande.
  \item Rapport de dépendance élevé : une forte proportion de jeunes limite l'épargne et l'investissement par tête.
  \item Pression sur les ressources et l'environnement : terres, eau, déforestation, urbanisation non planifiée.
\end{enumerate}

\paragraph{II. Une croissance qui peut être un atout (thèse « populationniste », Boserup).}
\begin{enumerate}[label=\Alph*.]
  \item Un marché intérieur en expansion et une main-d'{\oe}uvre abondante.
  \item La pression démographique comme stimulant de l'innovation (intensification agricole).
  \item Le dividende démographique : la baisse de la fécondité augmente la part des actifs (exemple de l'Asie de l'Est).
\end{enumerate}

\paragraph{III. Tout dépend des politiques publiques.}
\begin{enumerate}[label=\Alph*.]
  \item Accélérer la transition : santé de la reproduction, planification familiale, scolarisation des filles.
  \item Investir dans le capital humain et l'emploi des jeunes.
  \item Intégrer les variables démographiques dans la planification (rôle des démographes, des recensements, des données).
\end{enumerate}

\paragraph{Conclusion.}
La croissance démographique n'est ni un frein ni un moteur en soi : son effet dépend de la structure
par âge et de la capacité des États à investir dans l'éducation, la santé et l'emploi.
\emph{Ouverture} : le rôle de l'intégration régionale et des politiques de population dans l'Agenda 2063.

\end{document}
